% claim-ref: pothos#bud-growth-time
% claim-ref: pothos#established-plant-substrate-recipe
% claim-ref: pothos#groundcover-height
% claim-ref: pothos#indoor-flowering
% claim-ref: pothos#indoor-light-duration
% claim-ref: pothos#indoor-light-guidance
% claim-ref: pothos#lifespan
% claim-ref: pothos#preferred-temperature
% claim-ref: pothos#primary-method
% claim-ref: pothos#rooting-time
% claim-ref: pothos#seed-germination
% claim-ref: pothos#substrate-chemistry
% claim-ref: pothos#time-to-flowering
% claim-ref: pothos#transplant-readiness
% claim-ref: pothos#vine-length-ncsu
% claim-ref: pothos#vine-length-wisconsin
% claim-ref: pothos#watering-cue
% claim-ref: pothos#watering-interval
% claim-ref: pacifica-ca-coastal-context-v1#zip-94044-majority-hardiness-zone
\CompanionHeader{NUMBERS \& PACIFICA}{Pothos}{Provisional Epipremnum aureum}{Indoors; specimen identity, cultivar, variegation and household conditions remain unverified.}
\Context{Pacifica, coastal SF Bay Area: USDA 2023 ZIP 94044 majority is 10a (30-35 F). This is regional cold context, never a yard measurement or an indoor condition. [CLIMATE-USDA-PHZM-ZIP-94044]}
\CardRow{Light}{Filtered or bright indirect light. Low light can reduce variegation; avoid substantial direct sun. Indoor light duration is unknown: this category cannot be converted into exact sun hours. \textbf{[1,3]}}{Temperature}{Preferred published range: 70-90 F (21-32 C; C rounded conversion). This is species guidance, not a measured room range or a sharp survival boundary. \textbf{[3]}}
\CardRow{Water \& moisture}{Allow slight dry-down; water when the soil surface is dry and drain excess. A fixed interval is unknown because pot, medium, light and room conditions control dry-down. \textbf{[1,3]}}{Chemistry}{Numeric substrate pH unknown: the reviewed extension sources provide no applicable range. Use a well-drained, aerated container medium; do not infer chemistry from the city or a label. \textit{Proposed.} \textbf{[1,3]}}
\CardRow{Seed germination}{No practical, reproducible seed protocol established. Cultivated houseplants usually remain juvenile and do not flower; a germination interval is unknown. Node-bearing cuttings are the useful home method. \textbf{[1,3,4]}}{Vegetative establishment}{In warm conditions, buds can grow in 1-2 weeks and vines root in water or vermiculite in 3-4 weeks. These are separate milestones. Transfer after several roots and sustained growth. \textit{Proposed.} \textbf{[3,4]}}
\CardRow{Lifetime \& maturity}{No universal lifespan for a clonally renewed vine. Household first-flowering age is unknown; juvenile plants normally do not flower. Lack of flowers does not establish the plant's age. \textbf{[1,3]}}{Size \& growth}{Groundcover height: 6-8 in. Reported vine maxima differ: NCSU 40 ft; Wisconsin 65 ft. Retain their different species contexts; neither is a household container forecast. \textbf{[1,3]}}
\Panel{Established-plant recipe / 10 litres}{Proposed starting point for established plants; not propagation medium or the current mix.\par\vspace{3pt}\begin{tabular}{@{}p{.72\linewidth}rr@{}}\textbf{Ingredient / purpose} & \textbf{Volume} & \textbf{10 L}\\ commercial well-drained houseplant potting medium (base medium with water and nutrient retention) & 75\% & 7.5 L \\horticultural perlite (increase air space and drainage) & 25\% & 2.5 L \\\textbf{Total} & 100\% & 10 L\\\end{tabular}\par Observe drainage and dry-down before adjusting. \textit{Proposed.} \textbf{[1,3]}}
\Panel{Observe before adjusting}{Use observed indoor dry-down; outdoor ZIP zones do not measure the room. \textbf{[1,3]}\par Measured exposure / tank conditions: \hrulefill\par Observed establishment: \hrulefill\quad Date: \hrulefill}
\Sources{Evidence: 1 = POT-NCSU NC State; 3 = POT-WISC University of Wisconsin--Madison Extension; 4 = POT-NCSU-PROP NC State. Proposed = proposed starting point, not a published ratio or household measurement. Titles/URLs: entry and binder/contexts sources.yaml. Scopes: worksheets. Conversions rounded; source units authoritative.}{pothos / numbers}
